% =============================================================================
% Cover Letter — Solar Energy (Elsevier)
% Title: Spectral Co-Design of Agrivoltaic Modules for Multifunctional Net
%        Energy Benefit: A Quantum-Guided Digital Twin
% Date: November 6, 2026 (V2611 package)
% =============================================================================

\documentclass[10pt]{letter}
\usepackage[utf8]{inputenc}
\usepackage[T1]{fontenc}
\usepackage[hidelinks]{hyperref}
\usepackage[margin=1.6cm,bottom=1.0cm,top=1.4cm]{geometry}
\usepackage{siunitx}
\usepackage[version=4]{mhchem}
\usepackage{enumitem}
\setlist{nosep,leftmargin=1.2em,itemsep=1pt}
\DeclareSIUnit{\yr}{yr}
\DeclareSIUnit{\USD}{USD}
\DeclareSIUnit{\MWh}{MWh}
\DeclareSIUnit{\cubed}{^{3}}

\address{
Steve Cabrel Teguia Kouam\\
Department of Physics, Faculty of Science\\
University of Douala, Cameroon\\
\texttt{steve.teguia@facsciences-uy1.cm}
}

\begin{document}
\sloppy

\begin{letter}{Editorial Office\\
\textit{Solar Energy}\\
Elsevier}

\opening{Dear Editor,}

We submit our manuscript \textbf{``Spectral Co-Design of Agrivoltaic Modules for Multifunctional Net Energy Benefit: A Quantum-Guided Digital Twin''} for consideration as a Research Article in \textit{Solar Energy}. It is not under consideration elsewhere. The paper is about what a semi-transparent photovoltaic module should absorb, what it should hand to the crop below, and what that allocation is worth in water, carbon, and money --- solar energy conversion and spectral management at the heart of the journal's scope.

\noindent\textbf{Central question and significance.}
Can spectral co-design guided by non-Markovian quantum dynamics deliver a positive multifunctional net energy benefit (power $+$ water $+$ carbon $+$ diagnostics) in a field-constrained agrivoltaic stack, and which barriers remain physical rather than economic? Static semi-transparent modules resolve the food--energy conflict only by uniformly sacrificing output or yield. Here an organic photovoltaic layer coupled to a nanoparticle-on-mirror (NPoM) cavity shapes the transmitted spectrum actively: the quantum channel buys its diagnostics out of the near-infrared that even plant-transparent designs divert or discard, and the trade-off is quantified rather than asserted.

\noindent\textbf{What the analysis shows (all modeled, fully traceable).}
\begin{itemize}
\item The design law: a \qty{91.8}{\percent} local transport penalty, confined to a \qty{1}{\percent} sentinel area, costs \qty{0.9}{\percent} of canopy-scale device yield ($0.980 \to \num{0.971}$) and adds single-complex SERS diagnostics at no extra land.
\item The dual-band passbands at \qtylist{750;820}{\nm} lie outside PAR, so the photobiology and the crop-light budget do not compete; the land-use ledger returns a land-equivalent ratio of $\approx 1.4$.
\item The shield cuts crop evapotranspiration by \qty{28}{\percent} (\qty{460}{\L\per\m\squared\per\yr}) and yields a Net Ecological Benefit of \qty{72.6}{\kg\of{\ce{CO2e}}\per\m\squared\per\yr}; for a \qty{500}{\m\squared} cooperative module in Cameroon, payback is \qty{4.32}{\yr} at \qty{30}{\percent} subsidy (\qty{2.68}{\yr} with local materials).
\item Durability is the leading coefficient, not a footnote: reported outdoor organic-PV lifetimes bound a five-year replacement cycle, re-run through the carbon and payback accounts.
\end{itemize}

\noindent\textbf{Why \textit{Solar Energy}?}
The quantum methods appear only as the design engine for the module's spectral architecture; every claim that reaches a decision is optical, agronomic, or economic, and no revenue line in the techno-economics depends on the quantum gain. The paper is deliberately falsifiable: its exposure to the live debate over electronic coherence is quantified (all carried numbers survive the recent reattribution of long-lived signals to vibronic channels), and the Discussion names the three experiments that would change its conclusions.

\noindent\textbf{Novelty and positioning.}
To our knowledge, this is the first agrivoltaic stack whose spectral architecture is co-designed with numerically exact open-quantum-system dynamics, the first quantification of the local-to-global transport trade-off that makes a plasmonic diagnostic layer viable at canopy scale, and the first multifunctional TEA/LCA internalizing power, water, carbon, and sensing in a single account. Economic barriers are shown to be tractable; physical barriers (sub-\qty{2}{\nano\meter} cavity gaps, sentinel fouling, encapsulation lifetime) are stated honestly as the binding constraints.

\noindent\textbf{Compliance.}
The manuscript has been language-edited; all authors approve the submission and declare no competing interests. An earlier version was considered at \textit{Nature Energy}; we submit here as the more suitable venue. Subscription route (no APC). Corresponding author based in Cameroon.

\noindent\textbf{Suggested reviewers:} A.~Olaya-Castro (UCL), J.~Yuen-Zhou (UCSD), J.~J.~Baumberg (Cambridge), G.~A.~Barron-Gafford (Arizona)---covering quantum biology, polaritonic chemistry, NPoM/SERS, and agrivoltaics.

\medskip
\vspace{-0.5em}
Sincerely,

Steve Cabrel Teguia Kouam (on behalf of all authors)

\vspace{-1.4em}
\textit{Enclosures:} Manuscript, Supplementary Material, Highlights, Graphical Abstract.

\end{letter}

\end{document}
